\section{Evaluación y recompensa: cómo convertir tareas abiertas en objetivos entrenables}

\subsection{Evaluación basada en ejecución (Execution-based Evaluation)}
El curso utiliza evaluación por ejecución de scripts, en lugar de calificación subjetiva humana. Una vez terminada la tarea, un script de verificación revisa archivos, el estado de la página, el contenido de la salida o variables del sistema, y así determina el grado de finalización.

\begin{importantbox}{Las tareas abiertas también admiten una evaluación computable}
Incluso si la tarea admite múltiples rutas, basta con que el estado del resultado sea verificable para poder construir una evaluación automática. Lo importante es traducir la «intención de la tarea» en «condiciones verificables», por ejemplo, la existencia de un archivo, el valor de un campo, el estado de un elemento de la página o el valor de retorno de un comando.
\end{importantbox}

\begin{figure}[H]
\centering
\includegraphics[width=0.88\textwidth]{frame-04-reward.jpg}
\caption{Video aproximadamente en 00:19:35: el expositor presenta la recompensa basada en ejecución y la lógica de calificación por finalización parcial}
\end{figure}

\subsection{Recompensa por finalización parcial y estratificación de dificultad}
El curso enfatiza que no debe usarse únicamente un resultado binario de éxito o fracaso. En tareas complejas, una puntuación de finalización parcial refleja el progreso con mayor granularidad y mejora de manera notable el problema de las señales de entrenamiento dispersas. En la práctica, esto debe combinarse con una estratificación de dificultad, para evitar que el modelo se quede estancado durante mucho tiempo en subtareas de nivel bajo.

\begin{knowledgebox}{Recomendaciones para el diseño de la función de recompensa}
\begin{itemize}
    \item Dividir el objetivo en hitos clave y establecer recompensas por segmentos;
    \item Aumentar la penalización de las acciones destructivas, para restringir las conductas riesgosas;
    \item Aplicar un tratamiento de decaimiento a las operaciones repetitivas sin efecto, para inhibir las trayectorias cíclicas;
    \item Gestionar el versionado de los scripts de evaluación, para garantizar la reproducibilidad.
\end{itemize}
\end{knowledgebox}

\subsection{Requisitos de robustez del sistema de evaluación}
El propio script de evaluación también puede ser «atacado» o eludido. El curso recomienda incorporar en la capa del script ramas para casos anómalos, verificaciones de límites y restricciones contra trampas, para evitar que el modelo aprenda estrategias oportunistas ajenas al objetivo de la tarea.

\begin{warningbox}{Reward Hacking es un riesgo real}
Si la recompensa solo cubre un estado local, el Agent puede aprender a «hacer trampa con la puntuación» en lugar de «completar la tarea». Por ejemplo, generando un falso estado de finalización mediante el caché de la interfaz, o aprovechando una vulnerabilidad del script para evadir pasos clave. El diseño de la evaluación debe preceder a la iteración del modelo.
\end{warningbox}

\subsection{Resumen de la sección}
El límite superior del entrenamiento del Agent está determinado en gran medida por el sistema de evaluación. La evaluación basada en ejecución, la recompensa parcial y la robustez de los scripts determinan en conjunto si el modelo puede progresar de manera estable en tareas abiertas.

